\section{随机过程}

\answer{RS1}
非终止前缀为 $0,H,HT,HTH$，相应剩余均值记为 $e_0,e_1,e_2,e_3$。HTH 后遇 H，末尾最长有效前缀是 H，而不是空串。首步方程为
\[
 e_0=2+e_1,\quad e_1=2+e_2,\quad
 e_2=1+\tfrac12e_3+\tfrac12e_0,\quad
 e_3=1+\tfrac12e_1.
\]
从前两式 $e_1=e_0-2,e_2=e_0-4$，最后一式为 $e_3=e_0/2$。代入第三式得 $e_0-4=1+3e_0/4$，故 $\boxed{e_0=20}$。每个四次区块等于 HTHT 的概率为 $1/16$，保证均值有限。

\answer{RS2}
平稳几何比值满足 $\tfrac12r^2-r+\tfrac14=0$，选 $(0,1)$ 中的根 $r=1-1/\sqrt2$，所以
\[
 \boxed{\Prob_\pi(X\ge2)=\left(1-\frac1{\sqrt2}\right)^2},\qquad
 \boxed{\E_0\tau_0^+=\frac1{1-r}=\sqrt2}.
\]
一次清空或空队列服务均算正时间返回，因此该均值可以小于2。

\answer{RS3}
两步后切换当且仅当恰有一步切换，故 $2\alpha(1-\alpha)=3/8$，得到
\[
 \boxed{\alpha=1/4\quad\text{或}\quad\alpha=3/4}.
\]
两个参数对应的两步转移矩阵完全相同。对相同初始分布，所有偶数时刻联合分布因此相同，不能据此辨识。增加一步转移的信息即可区分；两个模型分别偏向保持与交替。

\answer{RS4}
给未处理 A、B 分别赋权3、2，处理一个任务使权重的条件期望减少1，故总任务数的均值有限。其均值满足 $a=1+b,b=1+a/3$，所以 $a=3,b=2$。若二阶矩为 $u,v$，独立的两个 B 子树给出
\[
 u=1+2b+v+b^2=9+v,\qquad
 v=1+\tfrac23a+\tfrac13u=3+\tfrac13u.
\]
这些等式也可先用于截断到有限代的任务树：均值不超过3、2，二阶矩递推被18、9控制。单调收敛证明二阶矩有限，然后解得
\[
 \boxed{\E T=3,\quad\E T^2=18,\quad\Var T=9}.
\]

\answer{RH1}
由 R11，$\rho=1/2$ 对应 $r=2-\sqrt3$，从而
\[
 \boxed{\E T=2+2\sqrt3}.
\]
第 $t$ 步是否仍需执行由此前决定，与当步是否重置独立，故重置总次数 $K$ 满足 $\E K=(1/2)\E T=\boxed{1+\sqrt3}$。路径上总费用为 $T+2K$，所以
\[
 \boxed{\E C=4+4\sqrt3}.
\]

\answer{RH2}
回撤差距在0处有自环，不能使用两端吸收区间的时间公式。R16 中 $d=3,D=12$，故
\[
 \boxed{\E\tau=12,\qquad\Var\tau=92},\qquad
 \boxed{\Prob(M_\tau\ge4)=\left(\frac34\right)^4=\frac{81}{256}}.
\]
最高点每提高一格，都须先赢下一次边界距离为1和3的竞争。

\answer{RH3}
R13 的分母为 $D=6(1/4)+1/2=2$，从而
\[
 \boxed{h_3^+=\frac{3/4+1/2}{2}=\frac58},\qquad
 \boxed{h_3^-=\frac{3/4}{2}=\frac38}.
\]
若初始方向等概率但未观察，按两种条件分布平均，成功概率为 $\boxed{1/2}$。这个平均不使位置过程自动成为无记忆游走。

\answer{RH4}
前0、1、2、3次互异的概率依次为 $1,1,2/3,2/9$，故
\[
 \E N=1+1+\frac23+\frac29=\frac{26}{9}.
\]
每次平均费用为 $7/3$。由 R18 的过去可测继续事件分解，
\[
 \boxed{\E C=\frac73\frac{26}{9}=\frac{182}{27}}.
\]
首次重复以前恰有 $N-1$ 个不同类型，最后一抽不增加类型数，所以其均值为 $\boxed{17/9}$。

\answer{RC1}
支付在上、下两状态分别为0、$1/2$。精确复制解
\[
 2\Delta+B=0,\qquad\tfrac12\Delta+B=\tfrac12,
\]
得到 $\boxed{\Delta=-1/3,B=2/3,V=1/3}$。若 $\Delta\ge0$，覆盖下状态要求 $\tfrac12\Delta+B\ge1/2$，于是初始成本 $\Delta+B\ge1/2+\Delta/2\ge1/2$。现金持有 $1/2$ 即可覆盖两状态，因此 $\boxed{V^{\rm super}=1/2}$。

\answer{RC2}
首次 H 后，下一次 H 的预测概率为 $2/3$，继续的期望增益为 $1/3$；首次 T 后，该增益为 $-1/3$，应停止。因此先抛，H 后再抛、T 后停，初始价值为
\[
 \boxed{V=\tfrac12(1+\tfrac13)+\tfrac12(-1)=\frac16}.
\]
直接停止价值0，不如观察一次；第二步的选择已比较全部可行动作。

\answer{RC3}
替换后的净支付为 $7/8-X$，因此在 $X<7/16$ 时替换，等号任取。令 $a=7/16$，积分得到
\[
 \boxed{V=\int_0^a(7/8-x)\dd x+\int_a^1x\dd x
 =\frac12+\frac{49}{256}=\frac{177}{256}}.
\]
不能使用独立重抽时的固定继续价值；这里替换价值依赖已观察的 $X$。

\answer{RC4}
全部通过的排序指标为 $c/(1-p)=(2,8,4/3)$，故顺序为 $(3,1,2)$，成本
\[
 \boxed{1+\tfrac14\cdot1+\tfrac14\cdot\tfrac12\cdot2=\frac32}.
\]
至少一项通过的指标为 $c/p=(2,8/3,4)$，故顺序 $(1,2,3)$，成本
\[
 \boxed{1+\tfrac12\cdot2+\tfrac12\cdot\tfrac14\cdot1=\frac{17}{8}}.
\]
每个后续成本都乘此前测试尚未确定结论的概率。

\answer{RG1}
$1,W_1,W_1^2-1$ 在平方可积空间中相互正交，后两者的方差为1和2。R25 给出 $\E I=1$、$\Cov(I,W_1)=1/2$、$\Var I=2e-5$。正态指数矩求导还给出
\[
 \E[e^{W_t-t/2}(W_1^2-1)]=t^2,\qquad
 \Cov(I,W_1^2-1)=\int_0^1t^2\dd t=\frac13.
\]
因此最佳预测及剩余误差为
\[
 \boxed{\widehat I=1+\frac12W_1+\frac16(W_1^2-1)}.
\]
\[
 \boxed{\E(I-\widehat I)^2=2e-5-\frac14-\frac1{18}=2e-\frac{191}{36}}.
\]
二次项使误差进一步减少 $1/18$，但有限阶多项式仍不等于完整条件均值。

\answer{RG2}
将 $T=2,b=0,a=1$ 代入 R27。时刻1的均值为 $6/8=3/4$，方差为 $1-1/2-3/8=1/8$，所以
\[
 \boxed{W_1\mid(W_2=0,A=1)\sim N(3/4,1/8)}.
\]
两侧点的条件协方差为
\[
 \boxed{\frac12-\frac{(1/2)(3/2)}2
 -\frac{3(1/2)(3/2)(3/2)(1/2)}8=-\frac{11}{128}}.
\]
面积约束使两处的偏离相互补偿。

\answer{RG3}
由条件反射公式，$T=2,b=0$ 时最大值 $M$ 的尾概率为 $\Prob(M\ge a)=e^{-a^2}$，$a\ge0$。因此
\[
 \boxed{\Prob(M\ge1\mid W_2=0)=e^{-1}},\qquad
 \boxed{\E(M\mid W_2=0)=\int_0^\infty e^{-a^2}\dd a=\frac{\sqrt\pi}{2}}.
\]

\answer{RG4}
R28 中 $k=\sqrt{2\lambda}=1$，故
\[
 \boxed{u(1)=\frac{\sinh1}{\sinh2}=\frac1{2\cosh1}},\qquad
 \boxed{\E(\tau\mid\text{成功})=2\coth2-\coth1=\tanh1}.
\]
最后一步用 $\sinh2=2\sinh1\cosh1$ 和 $\cosh2=\cosh^21+\sinh^21$。没有失效时钟时，成功概率为 $1/2$，成功条件时间均值为 $(4-1)/3=1$。失效同时降低成功率，并让成功样本偏向更快的路径。

\par\Needspace{8\baselineskip}
\subsection*{小节练习}
\answer{OR01}
把抛掷分成互不重叠的长度为 $n$ 的区块。每块全为正面的概率为 $2^{-n}$，且各块独立，故 $T_n$ 被 $n$ 倍的一个几何随机变量控制，因而可积。

令 $e_k$ 表示末尾已有 $k$ 次连续正面时，还需抛掷次数的均值，其中 $0\le k\le n$。终止边界为 $e_n=0$。出现反面会使连续正面数归零，因此
\[
 e_k=1+\tfrac12 e_{k+1}+\tfrac12e_0,
 \qquad 0\le k<n.
\]
从 $e_n=0$ 向后代入，得到
\[
 e_0=\left(1+\frac{e_0}{2}\right)
       \sum_{j=0}^{n-1}2^{-j}
     =\left(1+\frac{e_0}{2}\right)(2-2^{1-n}).
\]
移项后 $2^{-n}e_0=2-2^{1-n}$，故
\[
 \boxed{\E T_n=2^{n+1}-2}.
\]
等待下一次成功与等待一段连续成功不同；反面会清除已积累的进度。

\answer{OR07}
先设 $0<p<1$，记 $q=1-p$。以每次出发时所在地点的雨伞数 $X_t$ 为状态，状态空间为 $\{0,1,\ldots,n\}$。由于下一次从另一地点出发，转移为
\[
 P(0,n)=1,\qquad
 P(i,n-i)=q,\quad P(i,n-i+1)=p\quad(i\ge1).
\]
这条有限链不可约：从 $i<n$ 经一次不带伞、一次带伞可到 $i+1$；从 $i\ge2$ 经一次带伞、一次不带伞可到 $i-1$，而从1还可经 $n$ 到0。这些路径的概率均为正。

候选平稳分布为
\[
 \pi_0=\frac q{n+q},\qquad
 \pi_i=\frac1{n+q}\quad(1\le i\le n).
\]
流入0的概率为 $q\pi_n=\pi_0$。对 $1\le j<n$，流入概率为 $q\pi_{n-j}+p\pi_{n-j+1}=\pi_j$；流入 $n$ 的概率为 $\pi_0+p\pi_1=\pi_n$。这些概率和为1，故确为唯一平稳分布。

只有处于状态0且当次下雨时才会淋雨。有限不可约链的长期平均奖励等于平稳平均奖励，因此所求频率几乎必然为
\[
 \boxed{p\pi_0=\frac{p(1-p)}{n+1-p}}.
\]
该极限与初始的 $k$ 无关。$p=0$ 时从不淋雨；$p=1$ 时至多在第一次出行淋雨，之后每次出发地均有伞，长期频率也为0。上式在这两个端点仍然成立。

\answer{OR11}
记 $q=1-p$，并令 $u=u_1$。向左每次只走1，因此从 $n$ 到0必须依次经过 $n-1,\ldots,1$。在各下降首达时刻应用强马尔可夫性和平移不变性，得
\[
 u_n=u^n.
\]
这些下降之间允许出现任意多次向右移动，并不要求连续向左。对从1出发的第一步分类，有
\[
 u=p+qu_3=p+qu^3,
\]
即
\[
 (u-1)(qu^2+qu-p)=0.
\]
二次因子的正根为
\[
 r=\frac{\sqrt{1+4p/q}-1}{2}.
\]
若 $p\ge2/3$，则 $r\ge1$，而二次因子的另一根为负，故 $[0,1]$ 内唯一可用的根是1。

若 $p<2/3$，每步的均值为 $2-3p>0$，大数定律表明游走几乎必然趋于正无穷。此时 $u$ 不可能等于1：否则从任意整数到其左邻点的命中概率都是1，反复应用强马尔可夫性便会几乎必然访问所有负整数，与趋于正无穷矛盾。因此必须取 $u=r<1$。综上，
\[
 \boxed{u_n=
 \begin{cases}
 \left(\dfrac{\sqrt{1+4p/(1-p)}-1}{2}\right)^n,
       &0<p<2/3,\\[4pt]
 1,&2/3\le p<1.
 \end{cases}}
\]
首步方程可能有多个概率范围内的根，漂移或最小非负解的性质才决定应选哪一个。

\answer{OR16}
从任意未吸收状态，接下来的 $a+b$ 步若全向右便会吸收，其概率为 $2^{-(a+b)}$。因此 $\tau$ 有几何尾，所有阶矩均有限。

对称的独立增量使下列过程成为鞅：
\[
 S_n,\qquad S_n^2-n,\qquad S_n^3-3nS_n,
\]
\[
 S_n^4-6nS_n^2+3n^2+2n.
\]
例如分别展开 $(S_n\pm1)^j$ 再取平均即可核实。先在有界停时 $\tau\wedge N$ 应用可选停止，再令 $N\to\infty$。停止后的位置绝对值不超过 $\max(a,b)$，且 $\tau$ 的矩有限，故支配收敛允许这一极限。

设 $r=\Prob(S_\tau=b)$。前两个鞅分别给出
\[
 br-a(1-r)=0,\qquad
 \E\tau=\E S_\tau^2=b^2r+a^2(1-r),
\]
所以
\[
 \boxed{r=\frac a{a+b}},\qquad \boxed{\E\tau=ab}.
\]
为求二阶矩，令 $A=\E[\tau\ind_{\{S_\tau=b\}}]$、$B=\E[\tau\ind_{\{S_\tau=-a\}}]$。三次鞅给出
\[
 A+B=ab,\qquad
 bA-aB=\frac{\E S_\tau^3}{3}=\frac{ab(b-a)}3.
\]
解这两个线性方程，得
\[
 A=\frac{ab(2a+b)}{3(a+b)},\qquad
 B=\frac{ab(a+2b)}{3(a+b)}.
\]
于是
\[
 \E(\tau S_\tau^2)=b^2A+a^2B
       =\frac{ab(a^2+ab+b^2)}3,
\]
\[
 \E S_\tau^4=ab(a^2-ab+b^2).
\]
代入四次鞅的停止等式
\[
 0=\E S_\tau^4-6\E(\tau S_\tau^2)
       +3\E\tau^2+2\E\tau,
\]
最终得到
\[
 \boxed{\E\tau^2=\frac{ab(a^2+3ab+b^2-2)}3}.
\]

\answer{OR20}
令 $V(a,b)$ 表示剩余 $a$ 张红牌、$b$ 张黑牌时，今后可获得的最大期望增量。立即停止的增量是0，边界为
\[
 V(a,0)=a,\qquad V(0,b)=0.
\]
当 $a,b>0$，继续抽一张的期望价值为
\[
 C(a,b)=\frac{a[1+V(a-1,b)]+b[-1+V(a,b-1)]}{a+b},
\]
因此贝尔曼方程为
\[
 \boxed{V(a,b)=\max\{0,C(a,b)\}}.
\]
按 $a+b$ 从小到大计算即可得到全部状态的值。归纳证明最优性：在后继状态已经采用最优规则时，继续的价值恰为 $C(a,b)$；与停止的0比较就是当前能达到的最大值。牌数每次减少，故不存在无限期限或停时可积性问题。

从26红26黑出发，当前累计得分为 $s=b-a$，剩余牌数为 $N=a+b$。以下给出一条最优规则：$N=0$ 时结束；$N>0$ 且 $s\le0$ 时继续。对于正得分，按下表继续，其余情况停止：
\[
 \begin{array}{c|ccccc}
 s&1&2&3&4&5\\\hline
 \text{继续条件}&N\ge5&N\ge10&N\ge19&N\ge30&N\ge43
 \end{array}
\]
$s\ge6$ 时停止。表中只考虑实际可达的状态，因此 $N$ 与 $s$ 同奇偶。$s=1,N=3$ 时 $C=0$，停止或继续均最优，表中选择停止。

用上述递推进行精确分数计算，得到
\[
 \boxed{V(26,26)=\frac{41984711742427}{15997372030584}
       \approx2.624475549}.
\]
\par\smallskip\noindent\begin{minipage}{\linewidth}
递推也给出各状态的决策比较。以下程序计算初始精确值，避免引入舍入误差：
\begin{verbatim}
from fractions import Fraction as F
V = {}
for a in range(27):
    for b in range(27):
        if a == 0: V[a,b] = F(0)
        elif b == 0: V[a,b] = F(a)
        else:
            c = (a*(1+V[a-1,b])
                 + b*(-1+V[a,b-1]))/(a+b)
            V[a,b] = max(F(0), c)
print(V[26,26])
\end{verbatim}
\end{minipage}\par\smallskip
该问题的优势来自不放回抽样所提供的剩余牌信息；不能将每一步都当作胜率固定为 $1/2$ 的独立赌博。

\answer{OR22}
令 $V$ 为开始免费投掷前的最优价值。若当前点数为 $j$，接受的价值为 $j$；重新开始需先付1元，继续价值为 $V-1$。候选贝尔曼方程为
\[
 V=\frac16\sum_{j=1}^6\max\{j,V-1\}.
\]
$V=4$ 满足该式，因为
\[
 \frac{3+3+3+4+5+6}{6}=4.
\]
下面验证它确为全体策略的上界。令 $V_m$ 为最多投掷 $m$ 次的价值。$V_1=7/2\le4$，并且
\[
 V_{m+1}=\frac16\sum_{j=1}^6\max\{j,V_m-1\}.
\]
右边关于 $V_m$ 单调，故归纳得到 $V_m\le4$。若某策略在第 $N$ 次投掷停止且 $\E N<\infty$，强制它至迟在第 $m$ 次停止；所得趋于原所得，点数有界且费用被 $N$ 控制，故由支配收敛，其期望仍不超过4。若 $\E N=\infty$，由于最后点数至多6，期望净所得为负无穷，不可能更优。

因此最优决策为：1、2点重掷；4、5、6点接受；3点时接受或重掷均可。例如采用“首次出现至少3点即停”，最终点数在3、4、5、6上均匀，均值为 $9/2$；投掷次数为成功概率 $2/3$ 的几何变量，平均重掷费用为 $3/2-1=1/2$。于是
\[
 \boxed{\text{最大期望所得}=9/2-1/2=4}.
\]
这条策略达到上界且几乎必然结束，因而最优。

\answer{OR25}
对 $tW_t$ 用积分分部公式，得到
\[
 TW_T=\int_0^T W_t\,dt+\int_0^T t\,dW_t.
\]
因此
\[
 X=\int_0^T(T-t)\,dW_t,\qquad
 Y=\int_0^T t\,dW_t.
\]
两个被积函数都是确定的平方可积函数，所以两积分的均值均为0。伊藤等距及其交叉形式给出
\[
 \Var X=\int_0^T(T-t)^2dt=\frac{T^3}{3},\qquad
 \Var Y=\int_0^Tt^2dt=\frac{T^3}{3},
\]
\[
 \Cov(X,Y)=\int_0^Tt(T-t)dt=\frac{T^3}{6}.
\]
交叉形式也可将等距公式应用于两个被积函数的和，再展开平方得到。由于 $T>0$，两个方差均非零，故
\[
 \boxed{\Corr(X,Y)=\frac{T^3/6}{T^3/3}=\frac12}.
\]

\answer{OR26}
令
\[
 \tau=\inf\{t\ge0:B_t^{(1)}=-1\}.
\]
由反射原理，$\Prob(\tau\le t)=2\Phi(-1/\sqrt t)$，其中 $\Phi$ 为标准正态分布函数。因此 $\tau<\infty$ 几乎必然成立，其密度为
\[
 f_\tau(t)=\frac1{\sqrt{2\pi t^3}}e^{-1/(2t)},\qquad t>0.
\]
$\tau$ 只由第一坐标决定，与整个第二坐标过程独立，所以给定 $\tau=t$，有 $B_\tau^{(2)}\sim N(0,t)$。以 $Z=B_\tau^{(2)}$ 记第二坐标的增量。混合正态密度得到
\[
 f_Z(z)=\int_0^\infty\frac1{2\pi t^2}
            e^{-(1+z^2)/(2t)}dt
       =\frac1{\pi(1+z^2)}.
\]
这里被积函数非负，混合与积分交换由 Tonelli 定理保证；最后令 $u=1/t$ 即得积分值。于是 $Z$ 服从标准柯西分布，且
\[
 \boxed{\Prob(1+Z>0)
   =\int_{-1}^\infty\frac{dz}{\pi(1+z^2)}
   =\frac34}.
\]
独立的是停止时间与第二坐标的整个过程；随机时刻取值的无条件分布是上述混合分布，而不是固定方差的正态分布。

